\chapter{Options sur matières premières et couvertures structurées}\label{ch:m3:commodity-options-and-structured-hedges}

Chaque année, le ministère des Finances du Mexique achète des options de vente sur le pétrole qu'il prévoit d'exporter. En
décembre 2009, lorsque les cours s'étaient effondrés, ses options ont rapporté près de 5.1~milliards de dollars. En 2015, elles ont rapporté 6.4~milliards de dollars, et 2.7~milliards
l'année suivante. C'est l'une des plus grandes couvertures régulières de tout marché, et elle est construite
à partir des instruments de ce chapitre~: des options sur la moyenne d'un prix sur une période, achetées auprès
de banques qui les couvrent sur le marché des contrats à terme. Les producteurs, consommateurs et gouvernements de matières premières
se couvrent non pas sur un prix à une date, mais sur un revenu ou un coût sur plusieurs mois, et leurs instruments sont
adaptés en conséquence. Ce chapitre traite des \omterm{def:m3:commodity-options-and-structured-hedges:fo}{options sur contrats à terme} de matières premières, des \omterm{def:m3:commodity-options-and-structured-hedges:apo}{options sur prix moyen} et
des \omterm{def:m3:commodity-options-and-structured-hedges:cso}{options sur calendar spread}, les \omterm{def:m3:commodity-options-and-structured-hedges:swing}{contrats swing} du gaz et de l'électricité, les \omterm{def:m3:commodity-options-and-structured-hedges:collar}{collars} que vendent les producteurs
pour financer leur protection, ainsi que les \omterm{def:m3:commodity-options-and-structured-hedges:programme}{programmes de couverture} qui les combinent.

\section{Options sur contrats à terme de matières premières}

\begin{definition}[Option sur contrat à terme]\label{def:m3:commodity-options-and-structured-hedges:fo}
Une \emph{option sur contrat à terme}\index{option sur contrat à terme} est une option dont l'exercice confère une position dans un
contrat à terme (longue pour un call, courte pour un put) au prix d'exercice, la différence étant réglée
par la marge de variation du contrat à terme~; elle expire en général peu avant le contrat à terme sur lequel elle est écrite.
\end{definition}

Comme le sous-jacent est un contrat à terme, qui ne coûte rien à détenir, une \omterm{def:m3:commodity-options-and-structured-hedges:fo}{option sur contrat à terme} est évaluée sur le
forward, et non sur un spot majoré du carry~: la formule de Black pour les \omterm{def:m3:commodity-options-and-structured-hedges:fo}{options sur contrats à terme}, que dérive le livre~5,
dérive, utilise le cours du contrat à terme et actualise le paiement à l'échéance de l'option. Pour un contrat à terme valant 60~\$,
une option de vente à un an de strike 55~\$, avec une volatilité de 35~\% et un taux de 4~\%, vaut 5.51~\$ le baril. Au NYMEX,
l'option WTI cotée est américaine et expire à la clôture, trois jours ouvrés avant son contrat à terme~; son option WTI
sur prix moyen est européenne, réglée en espèces sur la moyenne mensuelle des prix de règlement du contrat à terme le plus proche
multipliée par 1\,000 barils, et expire le dernier jour ouvré du mois.

Les options sur matières premières diffèrent des options sur actions par la forme de leur sourire de volatilité (livre~1,
chapitre~25). Quand les chocs d'offre font monter les prix (panne de raffinerie, gel, guerre dans une zone productrice),
les calls hors de la monnaie peuvent coûter aussi cher que les puts, voire plus cher~; quand c'est un excédent qui est redouté, ce sont les puts.
Et les options à courte échéance portent la volatilité plus élevée des contrats à terme à courte échéance, l'\omterm{def:m3:forward-curves-storage-and-convenience-yield:samuelson}{effet Samuelson}
présenté au \cref{ch:m3:forward-curves-storage-and-convenience-yield}.

\section{Options sur prix moyen et options sur calendar spread}

Un producteur qui vend du pétrole chaque jour d'un mois est exposé au prix moyen du mois, et non à celui de son dernier
jour. Son option porte sur cette moyenne.

\begin{definition}[Option sur prix moyen]\label{def:m3:commodity-options-and-structured-hedges:apo}
Une \emph{option sur prix moyen}\index{option sur prix moyen} (average-price option, APO) verse à l'échéance, sur la moyenne arithmétique
d'un prix de référence sur une période (un mois, un trimestre, une année), par rapport à son strike~: un put
verse $\max(K - \bar P, 0)$ par unité. C'est le produit que les marchés de matières premières négocient sous ce nom~; le
livre~5, chapitre~16, évalue l'option asiatique générale.
\end{definition}

Une moyenne varie moins que sa dernière observation~: avec des fixings quotidiens, la moyenne d'un prix lognormal sur
la durée de vie de l'option présente environ $1/\sqrt 3$ de la volatilité du prix. Le put sur prix moyen de l'exemple
ci-dessus, sur douze fixings mensuels, vaut 2.66~\$ le baril par simulation, soit moins de la moitié du put vanille
de 5.51~\$ (\cref{fig:m3:commodity-options-and-structured-hedges:dist}).

\begin{omfigure}
\begin{tikzpicture}
\begin{axis}[width=11cm,height=4.8cm,ybar,bar width=4pt,
  xlabel={\$ par baril},ylabel={\% des trajectoires},
  tick label style={font=\small},label style={font=\small},
  xmin=10,xmax=150,ymin=0,ymax=22,grid=major,grid style={gray!25},
  legend columns=2,legend style={font=\footnotesize,at={(0.5,-0.33)},anchor=north,draw=none,fill=none,/tikz/every even column/.append style={column sep=8pt}},
  table/col sep=comma]
\addplot[fill=omDef!55,draw=omDef] table[x=mid,y=terminal]{figdata/markets-3/12-commodity-options-and-structured-hedges/dist.csv};
\addplot[fill=omThm!55,draw=omThm] table[x=mid,y=average]{figdata/markets-3/12-commodity-options-and-structured-hedges/dist.csv};
\legend{prix du contrat à terme dans un an,moyenne des douze fixings mensuels}
\end{axis}
\end{tikzpicture}

\omcaption{Le prix du contrat à terme en fin d'année et la moyenne de ses douze fixings mensuels, partant tous deux
de 60~\$ avec une volatilité de 35~\%, 200\,000 trajectoires simulées par intervalles de 5~\$~: la moyenne est bien moins
dispersée, si bien que les options qui s'y rapportent sont moins chères. Illustratif~; données~: tutoriel du chapitre.}\label{fig:m3:commodity-options-and-structured-hedges:dist}
\end{omfigure}

\begin{definition}[Option sur calendar spread]\label{def:m3:commodity-options-and-structured-hedges:cso}
Une \emph{option sur calendar spread}\index{option sur calendar spread} (calendar-spread option, CSO) est une option sur la différence entre
les prix de deux échéances de livraison d'une même matière première, par exemple un call versant $\max(F_{T_1} - F_{T_2}
- K, 0)$ sur le premier mois moins le second.
\end{definition}

Un propriétaire de capacités de stockage est naturellement acheteur d'\omterm{def:m3:commodity-options-and-structured-hedges:cso}{options sur calendar spread} (il y gagne lorsque les prix proches dépassent les prix éloignés
de plus que le coût de portage), et un raffineur ou un chargeur en est acheteur ou vendeur selon son
calendrier. Leurs prix dépendent de la corrélation entre les deux mois, qui est élevée mais baisse quand la forme de la
courbe change~; le livre~6, chapitre~16, modélise les options sur spread.

\section{Contrats swing}

Le gaz et l'électricité s'achètent sous des contrats qui laissent à l'acheteur une marge de flexibilité quant aux quantités prélevées, car
sa propre demande varie avec la météo.

\begin{definition}[Contrat swing, clause take-or-pay]\label{def:m3:commodity-options-and-structured-hedges:swing}
Un \emph{contrat swing}\index{contrat swing} (swing contract) permet à son détenteur de choisir chaque jour une quantité comprise entre un
minimum et un maximum autour d'une quantité contractuelle journalière, à un prix fixe ou indexé, sous réserve de limites sur
la quantité totale prélevée sur l'année. Une \emph{clause take-or-pay}\index{clause take-or-pay} (take-or-pay clause) oblige l'acheteur à
payer une quantité annuelle minimale, qu'il prélève ou non.
\end{definition}

\begin{omfigure}
\centering
\begin{tikzpicture}[font=\small]
\begin{axis}[width=11cm,height=4.4cm,
  xlabel={jour de l'année contractuelle (stylisé)},ylabel={quantité journalière},
  tick label style={font=\small},label style={font=\small},
  xmin=0,xmax=365,ymin=0,ymax=150,xtick=\empty,ytick={80,100,120},yticklabels={min,DCQ,max}]
\addplot[omDef,thick,dashed,domain=0:365] {120};
\addplot[omDef,thick,dashed,domain=0:365] {80};
\addplot[gray,thin,domain=0:365] {100};
\addplot[omThm,very thick,const plot] coordinates {(0,118) (60,120) (100,95) (160,80) (240,82) (300,112) (365,112)};
\node[font=\footnotesize,anchor=south west] at (axis cs:5,122) {prendre davantage quand le prix de marché est élevé};
\node[font=\footnotesize,anchor=north west] at (axis cs:125,78) {prendre le minimum quand il est bas};
\end{axis}
\end{tikzpicture}

\omcaption{Un \omterm{def:m3:commodity-options-and-structured-hedges:swing}{contrat swing}, schématique~: chaque jour, le détenteur notifie une quantité comprise entre un minimum et un maximum
autour de la quantité contractuelle journalière (DCQ)~; le total annuel doit rester dans des limites annuelles, et la
\omterm{def:m3:commodity-options-and-structured-hedges:swing}{clause take-or-pay} fait du total le plus bas un paiement dans tous les cas. Stylisé, sans données.}\label{fig:m3:commodity-options-and-structured-hedges:swing}
\end{omfigure}

Un \omterm{def:m3:commodity-options-and-structured-hedges:swing}{contrat swing} est une série d'options quotidiennes permettant d'acheter plus ou moins au prix du contrat, liées entre elles par les
limites annuelles~: prendre le maximum aujourd'hui consomme une flexibilité qui peut valoir davantage plus tard. Sa valeur
dépend donc de toute la trajectoire des prix, et son usage optimal relève d'une programmation dynamique,
comme pour la batterie étudiée au \cref{ch:m3:power-markets-trading}~; le modèle est celui du livre~6.

\begin{example}[Trois jours de swing]\label{ex:m3:commodity-options-and-structured-hedges:swing}
Un acheteur détient un contrat de trois jours à \qty{40}{EUR/MWh} avec une quantité contractuelle journalière de 100~MWh, un swing de
20~MWh dans les deux sens, et un total qui doit être égal à 300~MWh. Les prix de marché se révèlent être 30, 50 et 40. Prendre 80
le premier jour (en achetant les 20~MWh manquants sur le marché à 30) et 120 le deuxième (en revendant les 20 en trop à
50) vaut $20 \times (40 - 30) + 20 \times (50 - 40) = 400$ euros de plus que de prendre 100 chaque jour. Avec
des prix inconnus à l'avance, le détenteur doit décider chaque jour, et la flexibilité vaut moins que cette
valeur avec prévision parfaite.
\end{example}

\section{Collars et three-way collars}

Les puts coûtent de l'argent. Un producteur qui veut un plancher de prix mais refuse de le payer vend plutôt la hausse potentielle.

\begin{definition}[Collar et three-way collar]\label{def:m3:commodity-options-and-structured-hedges:collar}
Un \emph{collar}\index{collar} de producteur achète un put et vend un call, en général avec des strikes choisis de sorte que
les primes se compensent (un collar à coût nul)~: son prix réalisé reste compris entre les deux strikes. Un
\emph{three-way collar}\index{three-way collar} ajoute un put vendu sous le put acheté~: le producteur n'est
protégé qu'entre les deux strikes de put, et la prime supplémentaire lui permet de relever le plancher ou le cap.
\end{definition}

\begin{omfigure}
\begin{tikzpicture}
\begin{axis}[width=11cm,height=5.2cm,
  xlabel={prix moyen, \$/baril},ylabel={prix réalisé, \$/baril},
  tick label style={font=\small},label style={font=\small},
  xmin=20,xmax=110,ymin=20,ymax=110,ytick={40,60,80,100},grid=major,grid style={gray!25},
  legend columns=2,legend style={font=\footnotesize,at={(0.5,-0.25)},anchor=north,draw=none,fill=none,/tikz/every even column/.append style={column sep=8pt}},
  table/col sep=comma]
\addplot[gray,thick] table[x=x,y=unhedged]{figdata/markets-3/12-commodity-options-and-structured-hedges/payoffs.csv};
\addplot[omDef,very thick] table[x=x,y=collar]{figdata/markets-3/12-commodity-options-and-structured-hedges/payoffs.csv};
\addplot[omThm,very thick,dashed] table[x=x,y=threeway]{figdata/markets-3/12-commodity-options-and-structured-hedges/payoffs.csv};
\addplot[omProp,thick,dotted] table[x=x,y=put]{figdata/markets-3/12-commodity-options-and-structured-hedges/payoffs.csv};
\legend{non couvert,collar 50/75,three-way 60/45/70.24,put 55 net de prime}
\end{axis}
\end{tikzpicture}

\omcaption{Prix réalisé par baril d'un producteur, en regard du prix moyen de l'année~: le \omterm{def:m3:commodity-options-and-structured-hedges:collar}{collar} le maintient
entre 50~\$ et 75~\$~; le three-way ne protège qu'entre 60~\$ et 45~\$, puis suit le prix à la baisse avec
un coussin de 15~\$, et abandonne toute hausse au-delà de 70.24~\$~; le put acheté conserve tout le potentiel haussier au prix de
sa prime. Illustratif~; données~: le tutoriel du chapitre.}\label{fig:m3:commodity-options-and-structured-hedges:payoffs}
\end{omfigure}

\begin{example}[Coût de chaque structure]\label{ex:m3:commodity-options-and-structured-hedges:costs}
Avec un contrat à terme à 60~\$, une volatilité de 35~\% et une moyenne mensuelle sur un an, le \omterm{def:m3:commodity-options-and-structured-hedges:collar}{collar} 50/75 coûte
0.07~\$ le baril, presque nul~; le put sur prix moyen à 55 coûte 2.66~\$. Un three-way qui achète le put à 60 et
vend le put à 45 coûte le même 2.66~\$ s'il vend aussi un call à 70.24~\$. Un producteur disposant d'un budget de 50~\$
le baril serait, non couvert, inférieur de 8.43~\$ à ce budget dans les 5~\% des pires années~; avec le \omterm{def:m3:commodity-options-and-structured-hedges:collar}{collar} ou le
put, son prix réalisé des 5~\% pires années reste au-dessus du budget.
\end{example}

Le three-way a séduit les producteurs de pétrole parce qu'il paraît bon marché et protège contre des baisses modérées.
Sa faiblesse tient au put vendu~: en cas d'effondrement, la protection du producteur s'arrête au strike inférieur juste au moment
où elle est la plus nécessaire.

\begin{omfigure}
\begin{tikzpicture}
\begin{axis}[width=11cm,height=4.8cm,
  xlabel={prix réalisé sur l'année, \$/baril},ylabel={probabilité cumulée},
  tick label style={font=\small},label style={font=\small},
  xmin=25,xmax=110,ymin=0,ymax=1,grid=major,grid style={gray!25},
  legend columns=3,legend style={font=\footnotesize,at={(0.5,-0.3)},anchor=north,draw=none,fill=none,/tikz/every even column/.append style={column sep=8pt}},
  table/col sep=comma]
\addplot[gray,thick] table[x=unhedged,y=p]{figdata/markets-3/12-commodity-options-and-structured-hedges/revcdf.csv};
\addplot[omDef,very thick] table[x=collar,y=p]{figdata/markets-3/12-commodity-options-and-structured-hedges/revcdf.csv};
\addplot[omProp,very thick,dashed] table[x=put,y=p]{figdata/markets-3/12-commodity-options-and-structured-hedges/revcdf.csv};
\legend{non couvert,collar 50/75,put 55 net de prime}
\end{axis}
\end{tikzpicture}

\omcaption{Distribution du prix moyen réalisé d'un producteur sur une année, simulée à partir d'un contrat à terme à 60~\$ avec
une volatilité de 35~\%~: non couvert, avec le \omterm{def:m3:commodity-options-and-structured-hedges:collar}{collar} 50/75, et avec le put sur prix moyen à 55 net de sa prime de 2.66~\$.
Le \omterm{def:m3:commodity-options-and-structured-hedges:collar}{collar} tronque les deux queues~; le put ne tronque que la queue basse, au prix d'un coût. Illustratif~; données~: le
tutoriel du chapitre.}\label{fig:m3:commodity-options-and-structured-hedges:revcdf}
\end{omfigure}

\section{Programmes de couverture des producteurs, des consommateurs et des États}

\begin{definition}[Programme de couverture]\label{def:m3:commodity-options-and-structured-hedges:programme}
Un \emph{programme de couverture}\index{programme de couverture} est une politique permanente, fixée par un conseil d'administration ou un ministère,
qui fixe la part de la production ou de la consommation future couverte, sur quel horizon, avec quels
instruments, et la manière dont les couvertures sont exécutées et rapportées.
\end{definition}

Un producteur se couvre pour protéger son budget d'investissement et ses covenants d'endettement~; un consommateur (une compagnie aérienne, une
entreprise de services publics, un transformateur alimentaire) pour protéger ses marges~; un État pour protéger son budget. Ils choisissent
différemment~: les producteurs privilégient les swaps et les \omterm{def:m3:commodity-options-and-structured-hedges:collar}{collars}, qui ne coûtent rien d'emblée, les consommateurs les calls ou les \omterm{def:m3:commodity-options-and-structured-hedges:collar}{collars}, et les
États les puts, qui ont un coût de prime connu et laissent la hausse au Trésor public.

\begin{dated}{2026-09}{Couverture pétrolière du Mexique}\label{dat:m3:commodity-options-and-structured-hedges:mexico}
Le gouvernement mexicain achète chaque année des options de vente asiatiques (sur prix moyen) sur ses exportations de pétrole pour l'année
suivante, auprès de banques. Les versements déclarés s'élèvent à près de 5.1~milliards de dollars pour 2009, 6.4~milliards de dollars pour 2015 et 2.7
milliards de dollars pour 2016. Sur le coût, en janvier 2020, le ministre des Finances l'a estimé à environ 1~milliard de dollars par an, et un
vice-ministre l'a évalué à 1.2~milliards de dollars en moyenne sur la durée du programme, selon la presse.
\end{dated}

\section{Tutoriel~: valorisation d'un put sur prix moyen et d'un three-way}\label{tut:m3:commodity-options-and-structured-hedges:apo}

\begin{tutorial}
\textbf{Objectif.} Valoriser un put vanille sur un contrat à terme et un put sur prix moyen, et résoudre le strike du call du
d'un \omterm{def:m3:commodity-options-and-structured-hedges:collar}{three-way collar} dont le coût égale celui du put. \textbf{État final~:}
\cref{fig:m3:commodity-options-and-structured-hedges:dist,fig:m3:commodity-options-and-structured-hedges:payoffs}
et \cref{ex:m3:commodity-options-and-structured-hedges:costs}.
\begin{tutsteps}
\item \textbf{Les jambes.} Une jambe est un swap, un put ou un call, acheté ou vendu~; les structures sont des listes de jambes.
\omcode{code/firm/hedgeprog/firm_hedgeprog.py}{28}{52}{Jambes, collars et three-way collars.}\label{lst:m3:commodity-options-and-structured-hedges:legs}
\item \textbf{Simuler et valoriser.} Moyennes mensuelles d'un contrat à terme lognormal sans dérive, puis moyennes actualisées
des gains.
\omcode{code/firm/hedgeprog/firm_hedgeprog.py}{55}{68}{Moyennes par simulation et valeur d'une structure.}\label{lst:m3:commodity-options-and-structured-hedges:value}
\item \textbf{Exécuter} \texttt{m3\_options.vanilla\_vs\_average()}, \texttt{sovereign()},
\texttt{producer\_car()} et \texttt{fig\_options.py}.
\end{tutsteps}
\textbf{Pour aller plus loin.} Utiliser des fixings quotidiens plutôt que mensuels et comparer avec la règle en $1/\sqrt3$~;
ajouter un skew de volatilité (une volatilité plus élevée pour les bas strikes) et voir quelle structure devient la plus chère.
\end{tutorial}

\section{Construction~: l'évaluateur de programme de couverture}\label{bld:m3:commodity-options-and-structured-hedges:hedgeprog}

\begin{build}
\bfield{Rôle} La firme miniature structure des couvertures pour des producteurs et des consommateurs~: elle doit valoriser leurs structures, montrer leurs profils de gains
et mesurer ce que chacune fait aux pires années du client.
\bfield{Interface} \texttt{black76(f, k, t, sigma, r, call)}~; \texttt{Leg(kind, strike, side).payoff}~;
\texttt{collar}, \texttt{three\_way}~; \texttt{simulate\_averages(f0, sigma, t, steps, n, seed)}~;
\texttt{value(legs, settlements, r, t)}~; \texttt{hedged\_revenue(price, legs, premium)}~;
\texttt{cash\_flow\_at\_risk(revenue, budget, q)}.
\bfield{Règles} Gains par baril calculés sur le prix de règlement (la moyenne pour les jambes à prix moyen)~; simulation
graines fixes~; le contrat à terme est une martingale sous la mesure de valorisation~; le cash flow at risk est le déficit
du quantile inférieur par rapport au budget.
\bfield{Tests d'acceptation} \texttt{code/firm/hedgeprog/tests/}~: parité call-put de la formule de Black~; le put à prix moyen
moins cher que le put vanille, et la moyenne simulée centrée sur le terme~; un \omterm{def:m3:commodity-options-and-structured-hedges:collar}{collar}
qui borne le revenu~; la protection plafonnée d'un three-way.
\bfield{Pour aller plus loin} Sourires de volatilité selon le strike~; périodes de moyennage trimestrielles et annuelles~; les volumes du producteur
aussi incertains que ses prix.
\end{build}

\begin{omsources}
\item MoneyWeek, «~Mexico's lucrative Hacienda hedge~», 21 avril 2017.
\item Jain Family Institute, \emph{Mexico's Petroleum Hedging Program}, 2023.
\item Alto Nivel (avec Reuters), «~¿Qué son las coberturas petroleras y cuánto le cuestan a México?~», 10 janvier 2020.
\item F.~Black, «~The pricing of commodity contracts~», \emph{Journal of Financial Economics}, 1976.
\item CME Group, WTI average price options~; NYMEX Rulebook, chapitres 310 (Light Sweet Crude Oil Option) et 341 (WTI Average
Price Option) (copies archivées sur Internet Archive).
\end{omsources}

\section{Exercices}

\begin{exercise}[$\star$]\label{exo:m3:commodity-options-and-structured-hedges:1}
Un producteur vend un swap à 62~\$ pour une année de moyennes mensuelles. L'année, la moyenne s'établit à 55~\$. Calculer ce que paie
le swap et le prix réalisé du producteur.
\end{exercise}

\begin{exercise}[$\star$]\label{exo:m3:commodity-options-and-structured-hedges:2}
Pourquoi une \omterm{def:m3:commodity-options-and-structured-hedges:apo}{option sur prix moyen} est-elle moins chère qu'une option vanille de même strike et de même échéance~?
\end{exercise}

\begin{exercise}[$\star$]\label{exo:m3:commodity-options-and-structured-hedges:3}
Donner le prix réalisé par baril du \omterm{def:m3:commodity-options-and-structured-hedges:collar}{collar} 50/75 si la moyenne de l'année s'établit à 40~\$, 60~\$ et 90~\$.
\end{exercise}

\begin{exercise}[$\star\star$]\label{exo:m3:commodity-options-and-structured-hedges:4}
Pour le three-way 60/45/70.24, donner le prix réalisé si la moyenne de l'année s'établit à 30~\$, 50~\$, 65~\$ et 90~\$.
\end{exercise}

\begin{exercise}[$\star\star$]\label{exo:m3:commodity-options-and-structured-hedges:5}
Avec douze fixings mensuels, la volatilité du logarithme de la moyenne dans la simulation est de 21.3~\%. Comparer
cette valeur à $35\%/\sqrt3$ et expliquer l'écart.
\end{exercise}

\begin{exercise}[$\star\star$]\label{exo:m3:commodity-options-and-structured-hedges:6}
Pourquoi le propriétaire d'un stockage détient-il, en pratique, des \omterm{def:m3:commodity-options-and-structured-hedges:cso}{options sur calendar spread}~?
\end{exercise}

\begin{exercise}[$\star\star\star$]\label{exo:m3:commodity-options-and-structured-hedges:7}
\emph{Programmation.} Avec \texttt{black76}, vérifier la parité call-put pour l'\omterm{def:m3:commodity-options-and-structured-hedges:fo}{option sur contrat à terme} du texte, et donner
la valeur du call.
\end{exercise}

\begin{exercise}[$\star\star\star$]\label{exo:m3:commodity-options-and-structured-hedges:8}
\emph{Repérer l'erreur.} «~Notre \omterm{def:m3:commodity-options-and-structured-hedges:collar}{three-way collar} ne coûte rien et nous protège d'une chute sous 60~\$.~»
\end{exercise}

\section{Problème~: le put du ministère des Finances}

\begin{problem}[{Problème du week-end --- put annuel sur le pétrole d'un État}]\label{pb:m3:commodity-options-and-structured-hedges:1}
Un ministère des Finances prévoit d'exporter 250 millions de barils l'an prochain et les budgète au terme, à 60~\$. Il
peut acheter un put à prix moyen à 55~\$, ou consacrer la même prime à un three-way~: acheter le put 60, vendre le put 45
et vendre un call. La volatilité vaut 35~\%, le taux 4~\% et le moyennage est mensuel (à titre d'illustration).

\medskip\noindent\textbf{Partie I --- Le put.}
\begin{enumerate}
\item Quelle est la prime du put par baril~?
\item Combien coûte le programme, en dollars~?
\item Quelle part du revenu d'exportation budgété cela représente-t-il~?
\item Que rapporte le put si la moyenne de l'année s'établit à 45~\$~?
\item Et si elle s'établit à 30~\$~?
\end{enumerate}

\medskip\noindent\textbf{Partie II --- Le three-way.}
\begin{enumerate}[resume]
\item À quel strike le call doit-il être vendu pour la même prime~?
\item Que rapporte le three-way pour une moyenne de 45~\$~?
\item Et pour 30~\$~?
\item Combien le programme coûte-t-il au ministère si la moyenne de l'année s'établit à 90~\$~?
\item Quelle structure est la plus avantageuse selon les années~?
\end{enumerate}

\medskip\noindent\textbf{Partie III --- Exécution.}
\begin{enumerate}[resume]
\item Pourquoi acheter des options à prix moyen plutôt que des vanilles~?
\item Comment les banques vendeuses se couvrent-elles, et quel effet cela a-t-il sur le marché à terme~?
\item Pourquoi le ministère pourrait-il étaler ses achats sur plusieurs semaines et les garder secrets~?
\item Que fait le ministère si ses exportations sont inférieures à 250 millions de barils~?
\item Comment la prime doit-elle être comptabilisée dans le budget~?
\end{enumerate}

\medskip\noindent\textbf{Partie IV --- Jugement.}
\begin{enumerate}[resume]
\item Pourquoi un gouvernement préférerait-il des puts à des swaps~?
\item Qu'apprend l'historique des versements du programme sur sa valeur~?
\item Quand un three-way est-il un mauvais choix pour un État~?
\item Énoncer le \emph{résultat nommé}~: le coût du programme de puts en part du revenu, et ses versements pour
une année à 45~\$ et une année à 30~\$, comparés à ceux du three-way.
\item En une phrase~: qu'achète un État avec un put~?
\end{enumerate}
\end{problem}

\section{Questions d'entretien}

\begin{interviewq}[$\star$ \iqroles{trader}]\label{iq:m3:commodity-options-and-structured-hedges:1}
Qu'est-ce qu'une \omterm{def:m3:commodity-options-and-structured-hedges:apo}{option sur prix moyen}, et qui en a besoin~?
\end{interviewq}

\begin{interviewq}[$\star$ \iqroles{trader, banque}]\label{iq:m3:commodity-options-and-structured-hedges:2}
Expliquez à un directeur financier de producteur ce qu'est un \omterm{def:m3:commodity-options-and-structured-hedges:collar}{collar} à coût nul.
\end{interviewq}

\begin{interviewq}[$\star\star$ \iqroles{chercheur}]\label{iq:m3:commodity-options-and-structured-hedges:3}
Pourquoi la volatilité effective d'une moyenne vaut-elle environ $1/\sqrt3$ de celle du prix~?
\end{interviewq}

\begin{interviewq}[$\star\star$ \iqroles{banque, risque}]\label{iq:m3:commodity-options-and-structured-hedges:4}
Vous avez vendu un put à prix moyen important à un État. Comment le couvrez-vous au fil de l'année~?
\end{interviewq}

\begin{interviewq}[$\star\star$ \iqroles{trader}]\label{iq:m3:commodity-options-and-structured-hedges:5}
Qu'est-ce qu'un \omterm{def:m3:commodity-options-and-structured-hedges:swing}{contrat swing}, et pourquoi vaut-il plus qu'un contrat à volume fixe~?
\end{interviewq}

\begin{interviewq}[$\star\star\star$ \iqroles{développeur, chercheur}]\label{iq:m3:commodity-options-and-structured-hedges:6}
Concevez un outil qui propose et compare des structures de couverture pour un producteur, à partir de son budget et de ses covenants.
\end{interviewq}
